\documentclass[../DoAn.tex]{subfiles}
\begin{document}

\section{Mở đầu chương}

Chương này trình bày phương pháp đề xuất cho hệ thống trợ lý học vụ
\texttt{RAG\_v2}. Mục tiêu của hệ thống là trả lời câu hỏi của sinh viên dựa
trên các tài liệu học vụ có căn cứ, đồng thời hỗ trợ hội thoại nhiều lượt, truy
vết nguồn, cập nhật tài liệu quản trị và xử lý các câu hỏi phức tạp bằng
Agentic RAG.

Phương pháp được xây dựng trực tiếp từ kiến trúc mã nguồn hiện tại. Hệ thống
không chỉ là một pipeline RAG tuyến tính, mà là một nền tảng nhiều lớp gồm:

\begin{itemize}
    \item backend FastAPI cung cấp API chat, streaming, quản trị tài liệu, xác
    thực, phiên hội thoại và các tính năng mobile;
    \item tầng orchestration \texttt{RAGPipeline} điều phối classic RAG, smart
    routing, agent, streaming và fallback;
    \item tầng truy hồi lai dùng Qdrant, Elasticsearch, BGE-M3, E5 multilingual
    và BGE reranker;
    \item tầng hiểu truy vấn gồm complexity routing, domain classification,
    reflection, entity extraction và decomposition;
    \item tầng Agentic RAG dựa trên LangGraph cho các câu hỏi so sánh, nhiều
    nguồn hoặc cần công cụ;
    \item tầng dữ liệu, chunking, indexing, cache, logging, đánh giá và giao
    diện web/mobile.
\end{itemize}

Khác với cách thiết kế một chatbot chỉ gọi LLM, phương pháp đề xuất đặt trọng
tâm vào ba nguyên tắc: \textbf{có căn cứ}, \textbf{đúng miền dữ liệu} và
\textbf{có khả năng quan sát}. Mỗi câu trả lời cần được sinh từ context truy hồi,
có metadata để truy vết, có cơ chế tránh dùng sai văn bản và có fallback khi dữ
liệu nội bộ không đủ.

\section{Mục tiêu thiết kế}

\subsection{Yêu cầu chức năng}

Hệ thống được thiết kế để đáp ứng các nhóm yêu cầu chức năng sau:

\begin{itemize}
    \item \textbf{Hỏi đáp học vụ}: trả lời câu hỏi về chương trình đào tạo, quy
    định, kế hoạch học tập, thủ tục và hỗ trợ sinh viên.
    \item \textbf{Hội thoại nhiều lượt}: giữ lịch sử theo phiên, xử lý câu hỏi
    nối tiếp, tham chiếu như ``ngành của em'', ``môn đó'', ``so với khóa kia''.
    \item \textbf{Cá nhân hóa có kiểm soát}: dùng hồ sơ đăng nhập như mã ngành,
    khóa, email hoặc mã sinh viên khi câu hỏi hiện tại thật sự phụ thuộc vào
    người dùng; không để lịch sử hoặc hồ sơ làm nhiễu các câu hỏi chung.
    \item \textbf{Streaming}: hỗ trợ luồng trả lời thời gian thực qua SSE cho
    web và mobile.
    \item \textbf{Agentic reasoning}: xử lý câu hỏi cần so sánh, nhiều nguồn,
    nhiều bước hoặc cần hỏi lại người dùng.
    \item \textbf{Quản trị tài liệu}: cho phép admin upload, convert, clean,
    review, chunk, chọn chunk, index và rollback tài liệu.
    \item \textbf{Quan sát và đánh giá}: ghi lại route, nguồn, filter, timings,
    agent trace và hỗ trợ dashboard/evaluation.
\end{itemize}

\subsection{Yêu cầu phi chức năng}

Các yêu cầu phi chức năng chính gồm:

\begin{itemize}
    \item \textbf{Độ chính xác}: ưu tiên tài liệu đúng collection, đúng ngành,
    đúng khóa, đúng văn bản hiện hành và đúng thời điểm.
    \item \textbf{Khả năng mở rộng}: dữ liệu được tách theo collection; retrieval
    có thể chạy song song; Qdrant, Elasticsearch, MongoDB và Redis được tách
    thành các hạ tầng độc lập.
    \item \textbf{Hiệu năng}: dùng cache, candidate pool, rerank top-k, query-only
    cache khi an toàn và streaming để giảm cảm nhận độ trễ.
    \item \textbf{Tính bền vững vận hành}: cấu hình qua \texttt{Settings} và
    \texttt{.env}; các override LLM của admin có thể persist trong MongoDB.
    \item \textbf{An toàn dữ liệu}: Bearer JWT được ưu tiên hơn identity trong
    request body; admin route dùng RBAC; thông tin cá nhân không nên đi vào
    truy hồi nếu không cần.
\end{itemize}

\section{Kiến trúc tổng thể}

\subsection{Mô hình phân tầng}

Hệ thống được tổ chức theo kiến trúc phân tầng. Mỗi tầng có trách nhiệm rõ ràng
và giao tiếp với tầng khác qua contract cụ thể.

\begin{verbatim}
Nguoi dung Web/Mobile/Admin
        |
        v
FastAPI API Layer
  - /chat, /chat/v3, /chat/stream
  - /admin/documents, /retrieval/search, /metrics, /auth
        |
        v
RAGPipeline / DocumentPipeline
  - query_v3, query, query_agent, query_stream
  - convert, clean, chunk, embed_and_index, rollback
        |
        v
Query Understanding + Agent
  - ComplexityRouter, QueryRouter, QueryReflector
  - QueryDecomposer, ReActAgent, tool adapters
        |
        v
Retrieval Layer
  - RetrievalService
  - MultiCollectionSearch
  - Qdrant + Elasticsearch + BGE reranker
        |
        v
Data / Persistence / Cache
  - data collections, uploads
  - MongoDB, Redis
\end{verbatim}

Điểm trung tâm của thiết kế là \texttt{RAGPipeline}. FastAPI chỉ nhận request,
xác thực, ánh xạ user context và chuẩn hóa response; các quyết định nghiệp vụ
về route, truy hồi, agent, fallback và generation nằm trong pipeline.

\subsection{Các module chính trong source code}

\begin{table}[h]
\centering
\begin{tabular}{|l|p{10.2cm}|}
\hline
\textbf{Module} & \textbf{Vai trò chính} \\
\hline
\texttt{api/} & FastAPI app factory, lifespan, middleware, chat route, streaming,
admin upload route, metrics, retrieval diagnostic, session, lookup, bookmark,
feedback và notification. \\
\hline
\texttt{pipeline/} & Điều phối runtime chat/RAG qua \texttt{RAGPipeline} và pipeline
quản trị tài liệu qua \texttt{DocumentPipeline}. \\
\hline
\texttt{query/} & Complexity routing, domain classification, query reflection,
entity extraction, decomposition và structured query parsing. \\
\hline
\texttt{retrieval/} & Qdrant vector search, Elasticsearch keyword search, metadata
filtering, multi-collection fusion, validity filtering và reference resolution. \\
\hline
\texttt{agent/} & LangGraph ReAct agent, planner-executor, tool schema, tool
adapter, agent state và trace. \\
\hline
\texttt{chunking/} & Các chunker cho tài liệu pháp lý, CTDT, kế hoạch và STSV; làm
giàu metadata cho chunk. \\
\hline
\texttt{document\_loader/} & Convert PDF sang Markdown bằng Docling/PyMuPDF4LLM và
clean Markdown trước khi chunking. \\
\hline
\texttt{embedding/} & BGE-M3, E5 multilingual và interface embedding dùng trong
retrieval/classification/indexing. \\
\hline
\texttt{reranking/} & BGE reranker cross-encoder và factory tạo reranker. \\
\hline
\texttt{models/} & MongoDB access, user, session, turn, query log, agent trace,
document và system config. \\
\hline
\texttt{cache/} & Redis session store, history cache, LLM cache và sliding-window
rate limiter. \\
\hline
\texttt{frontend/}, \texttt{mobile/}, \texttt{packages/} & React web app, Expo
mobile app và shared TypeScript API/types/stores. \\
\hline
\texttt{evaluation/} & Offline current-policy eval, historical conversation eval,
benchmark retrieval và dashboard artifacts. \\
\hline
\end{tabular}
\caption{Bản đồ module chính của hệ thống}
\end{table}

\subsection{Stack công nghệ}

\begin{table}[h]
\centering
\begin{tabular}{|l|p{9.8cm}|}
\hline
\textbf{Thành phần} & \textbf{Công nghệ và implementation} \\
\hline
API backend & FastAPI, app factory/lifespan trong \texttt{api/main.py}, SSE cho
\texttt{/chat/stream}. \\
\hline
LLM chat/reflection & Provider qua \texttt{llm.create\_llm()}, mặc định cấu hình
Gemini; model đọc từ \texttt{Settings.chat\_model} và
\texttt{Settings.reflection\_model}. \\
\hline
Agent & LangGraph \texttt{StateGraph}, LangChain tool schema, local LM Studio
tool-calling model theo \texttt{agent\_model}, synthesis provider cấu hình riêng. \\
\hline
Vector store & Qdrant, mỗi collection có named vectors \texttt{bge\_m3} và
\texttt{e5}, mỗi vector 1024 chiều, cosine distance. \\
\hline
Keyword store & Elasticsearch, index theo collection, BM25 trên \texttt{text} và
\texttt{title}, kèm metadata filtering. \\
\hline
Embedding & \texttt{BAAI/bge-m3} và \texttt{intfloat/multilingual-e5-large}. \\
\hline
Reranker & \texttt{BAAI/bge-reranker-v2-m3} qua \texttt{BGEReranker}. \\
\hline
Persistence & MongoDB cho users, sessions, turns, query logs, agent traces,
documents, document chunks và system config. \\
\hline
Cache & Redis tùy chọn cho session, history, LLM cache và rate limiting. \\
\hline
Web client & React 18, Vite, React Router, TanStack Query, shadcn/Radix,
Tailwind CSS. \\
\hline
Mobile client & Expo/React Native, React Navigation, TanStack Query, SecureStore,
MMKV fallback, \texttt{react-native-sse}. \\
\hline
\end{tabular}
\caption{Stack công nghệ của phương pháp đề xuất}
\end{table}

\section{Nguồn dữ liệu và tổ chức tri thức}

\subsection{Các collection tri thức}

Tri thức của hệ thống được chia thành các collection theo miền nghiệp vụ thay
vì gom tất cả tài liệu vào một kho chung. Thiết kế này giúp routing, metadata
filtering và evaluation rõ ràng hơn.

\begin{table}[h]
\centering
\begin{tabular}{|l|p{4.0cm}|p{5.8cm}|}
\hline
\textbf{Collection} & \textbf{Nội dung} & \textbf{Metadata/logic quan trọng} \\
\hline
\texttt{ctdt} & Chương trình đào tạo, học phần, tín chỉ, học kỳ, tiên quyết. &
\texttt{major\_code}, \texttt{major\_name}, \texttt{course\_code},
\texttt{course\_name}, \texttt{semester}. \\
\hline
\texttt{quydinh} & Quy chế, học bổng, tốt nghiệp, ngoại ngữ, quy định học vụ. &
\texttt{applicable\_cohort}, \texttt{applicable\_major},
\texttt{effective\_date}, \texttt{document\_type}, lineage. \\
\hline
\texttt{kehoach} & Kế hoạch học kỳ, lịch đăng ký, thông báo, deadline. &
\texttt{date\_str}, freshness sorting, recency bonus, web fallback cho truy vấn
mới nhất. \\
\hline
\texttt{stsv} & Sổ tay sinh viên, biểu mẫu, thủ tục, hỗ trợ sinh viên. &
\texttt{doc\_id}, \texttt{title}, \texttt{type\_doc},
\texttt{section\_context}, thường không prefilter theo ngành/khóa. \\
\hline
\texttt{test} & Collection hợp lệ cho upload/dev. & Phụ thuộc metadata tài liệu
upload. \\
\hline
\end{tabular}
\caption{Các collection tri thức và vai trò runtime}
\end{table}

\subsection{Quy ước metadata}

Mỗi chunk sau khi xử lý cần có text/content, id ổn định, metadata nguồn, thứ tự
chunk và các trường miền dữ liệu. Metadata không chỉ phục vụ hiển thị nguồn mà
còn là điều kiện lọc trước khi vector search. Ví dụ:

\begin{itemize}
    \item truy vấn về CTDT ngành \texttt{IT-E6} ưu tiên lọc
    \texttt{major\_code=IT-E6};
    \item truy vấn về quy định khóa \texttt{K70} dùng
    \texttt{applicable\_cohort};
    \item truy vấn kế hoạch mới nhất dùng \texttt{date\_str} và sort theo ngày
    đăng;
    \item văn bản upload qua admin dùng \texttt{document\_id} để rollback và
    xóa dữ liệu đã index.
\end{itemize}

\subsection{Lineage và hiệu lực văn bản}

Với nhóm quy định, hệ thống sử dụng \texttt{data/document\_lineage.json} để mô
tả tài liệu còn hiệu lực hoặc đã bị thay thế. \texttt{ValidityFilter} dùng registry
này để loại các văn bản superseded khi có đủ kết quả thay thế. Nếu lọc làm số
kết quả còn lại quá ít, hệ thống giữ lại kết quả gốc để tránh tạo context rỗng.

\section{Pipeline xử lý và quản trị tài liệu}

\subsection{Luồng xử lý tài liệu ngoại tuyến}

Các nguồn dữ liệu ban đầu được xử lý bằng pipeline theo miền:

\begin{enumerate}
    \item thu thập hoặc crawl dữ liệu từ PDF, API hoặc website chính thức;
    \item convert PDF sang Markdown bằng Docling, PyMuPDF4LLM hoặc output OCR;
    \item làm sạch Markdown/text, chuẩn hóa heading, bảng, ký tự và HTML còn sót;
    \item chunking theo chiến lược phù hợp với loại tài liệu;
    \item enrich metadata về ngành, khóa, ngày, loại tài liệu và cấu trúc;
    \item embed bằng BGE-M3 và E5;
    \item index vào Qdrant và Elasticsearch.
\end{enumerate}

Các script vận hành trong \texttt{scripts/} như \texttt{index\_kehoach.py},
\texttt{index\_quydinh.py}, \texttt{index\_stsv.py}, \texttt{index\_to\_es.py}
và \texttt{auto\_crawler.py} hỗ trợ cập nhật dữ liệu độc lập với runtime API.

\subsection{Chunking theo loại tài liệu}

\begin{table}[h]
\centering
\begin{tabular}{|l|p{4.2cm}|p{5.7cm}|}
\hline
\textbf{Chunker} & \textbf{File chính} & \textbf{Chiến lược} \\
\hline
Recursive & \texttt{chunking/chunker/recursive\_chunker.py} & Chia theo heading
Markdown và recursive splitter, bảo vệ bảng, thêm section context. \\
\hline
Hierarchical legal & \texttt{hierarchical\_legal\_chunker.py},
\texttt{olmocr\_legal\_chunker.py} & Nhận diện chương, điều, khoản, phụ lục;
tạo parent/child chunk với hierarchy metadata. \\
\hline
\texttt{kehoach} & \texttt{chunking/chunker/kehoach\_chunker.py} & Chia bài kế
hoạch/thông báo, giữ tiêu đề, URL, ngày đăng và category. \\
\hline
\texttt{stsv} & \texttt{chunking/chunker/stsv\_chunker.py} & Chia nội dung JSON
theo phần Roman, mục số, văn xuôi và thêm context prefix. \\
\hline
\end{tabular}
\caption{Các chiến lược chunking theo source code hiện tại}
\end{table}

Điểm quan trọng của chunking là không chỉ tạo các đoạn text ngắn hơn, mà còn giữ
ngữ cảnh để đoạn truy hồi có thể đứng độc lập. Với tài liệu pháp lý, điều/khoản
được biểu diễn trong metadata. Với CTDT, bảng học phần cần được bảo vệ để không
tách mất header. Với STSV, prefix gồm tiêu đề, loại tài liệu và section giúp
chunk ngắn vẫn có ngữ cảnh.

\subsection{Admin DocumentPipeline}

Bên cạnh pipeline ngoại tuyến, hệ thống có luồng quản trị tài liệu trong
\texttt{DocumentPipeline}. Luồng này được gọi qua các endpoint \texttt{/admin}
và yêu cầu quyền admin.

\begin{verbatim}
uploaded
  -> converting -> converted
  -> cleaning   -> cleaned
  -> chunking   -> chunked
  -> embedding  -> indexed
\end{verbatim}

\texttt{DocumentPipeline} cung cấp các phương thức chính:
\texttt{convert\_pdf()}, \texttt{clean()}, \texttt{chunk()},
\texttt{embed\_and\_index()}, \texttt{run\_full\_pipeline()},
\texttt{delete\_indexed\_data()} và \texttt{rollback()}. Dữ liệu quản trị được
lưu ở Mongo collections \texttt{documents} và \texttt{document\_chunks}; nội dung
đã index được ghi vào Qdrant và Elasticsearch. Khi index, policy
\texttt{utils.chunk\_indexing.is\_indexable\_chunk()} giúp bỏ qua parent/header
chunk không phù hợp cho retrieval trực tiếp.

\section{Luồng runtime của RAGPipeline}

\subsection{Khởi tạo backend}

Entrypoint backend là \texttt{api/main.py}. Khi FastAPI khởi động, lifespan thực
hiện các bước:

\begin{enumerate}
    \item load \texttt{.env} và tạo \texttt{Settings};
    \item merge admin LLM override từ Mongo \texttt{system\_config/llm\_config}
    nếu có;
    \item khởi tạo \texttt{MongoLogger} khi MongoDB được bật;
    \item khởi tạo Redis manager, session store, history cache, LLM cache và
    rate limiter khi Redis được bật;
    \item tạo một \texttt{RAGPipeline} dùng chung \texttt{Settings};
    \item tạo Mongo indexes và warm up agent LLM nếu khả dụng;
    \item tùy chọn schedule auto crawler khi \texttt{crawler\_enabled=True}.
\end{enumerate}

Thiết kế này giúp toàn bộ backend dùng một cấu hình hiệu lực duy nhất. Admin có
thể cập nhật cấu hình LLM runtime qua API quản trị; pipeline chuẩn bị client mới,
persist override rồi commit thay đổi dưới runtime lock.

\subsection{Các entrypoint chat}

\texttt{RAGPipeline} có bốn entrypoint chính:

\begin{itemize}
    \item \texttt{query()}: classic RAG tuyến tính;
    \item \texttt{query\_v3()}: smart entrypoint dùng bởi web/mobile/API;
    \item \texttt{query\_agent()}: chạy agent trực tiếp;
    \item \texttt{query\_stream()}: retrieval trước rồi stream token qua
    \texttt{/chat/stream}.
\end{itemize}

\texttt{query\_v3()} là entrypoint quan trọng nhất. Nó quyết định câu hỏi đi theo
nhánh chitchat, classic RAG, decomposed RAG, agent hoặc clarify.

\subsection{Smart query flow}

\begin{verbatim}
query_v3(question, history, user_context)
  -> ComplexityRouter
     -> chitchat
        -> _handle_chitchat()
     -> complex + personal_check
        -> clarify, khong truy hoi khi thieu du lieu ca nhan
     -> simple
        -> query()
     -> complex + comparison/multi_source
        -> QueryDecomposer -> _query_decomposed()
     -> complex general
        -> query_agent()
           -> fallback query() neu agent loi va require_agent=False
\end{verbatim}

Các mode response thường gặp gồm \texttt{chitchat}, \texttt{clarify},
\texttt{rag\_v2}, \texttt{rag\_v2\_decomposed}, \texttt{agent} và
\texttt{rag\_v2\_fallback}. Với câu hỏi kiểm tra điều kiện cá nhân như tốt
nghiệp, nếu thiếu CPA/GPA, số tín chỉ, ngoại ngữ, GDTC, GDQP-AN hoặc trạng thái
kỷ luật, hệ thống trả \texttt{clarify} thay vì tự suy đoán từ tài liệu chung.

\subsection{Classic RAG flow}

Classic RAG được triển khai chủ yếu trong \texttt{pipeline/flows.py}. Luồng xử lý
chính gồm:

\begin{enumerate}
    \item trim lịch sử hội thoại theo ngân sách context;
    \item kiểm tra query-only cache khi truy vấn không mang tính dynamic;
    \item route domain bằng \texttt{QueryRouter};
    \item chọn collection bằng \texttt{CollectionSelector};
    \item rewrite truy vấn và extract entities bằng \texttt{QueryReflector};
    \item xây metadata filters theo collection;
    \item embed truy vấn bằng BGE-M3 và E5;
    \item chạy \texttt{MultiCollectionSearch};
    \item retry với chiến lược nới lỏng khi không có kết quả;
    \item deduplicate candidate;
    \item rerank bằng BGE reranker;
    \item áp \texttt{ValidityFilter};
    \item resolve dẫn chiếu pháp lý bằng \texttt{ReferenceResolver};
    \item format context, thêm profile note khi phù hợp;
    \item sinh câu trả lời bằng LLM;
    \item ghi cache, logging, timings và metadata;
    \item tùy chọn chạy self-eval/Tavily fallback ở non-streaming path.
\end{enumerate}

\subsection{Streaming flow}

\texttt{query\_stream()} phục vụ \texttt{/chat/stream}. Endpoint SSE gửi các event:

\begin{verbatim}
{"type":"session","session_id":"..."}
{"type":"token","delta":"..."}
{"type":"metadata", ...}
{"type":"done"}
\end{verbatim}

Streaming path vẫn chạy retrieval trước để có context grounded, sau đó mới phát
token. Khác với non-streaming path, streaming không chạy post-generation
self-eval/Tavily fallback để giữ semantics phát token ổn định.

\section{Hiểu truy vấn và định tuyến}

\subsection{ComplexityRouter}

\texttt{ComplexityRouter} là tầng đầu tiên trong \texttt{query\_v3()}. Router này
trả về \texttt{tier}, \texttt{reason}, \texttt{confidence} và có thể thêm
\texttt{complex\_subtype}. Các subtype phức tạp chính gồm:
\texttt{comparison}, \texttt{multi\_source}, \texttt{personal\_check} và
\texttt{general}.

Việc tách complexity trước retrieval giúp hệ thống tránh chi phí không cần thiết:
câu chào hỏi không truy hồi; câu hỏi đơn giản đi classic RAG; câu hỏi so sánh đi
decomposition hoặc agent.

\subsection{Domain classification và collection selection}

\texttt{DomainClassifier} có hai bước: phân loại intent
\texttt{chitchat}/\texttt{rag}/\texttt{tool\_search}, sau đó phân loại domain RAG
trong \texttt{ctdt}, \texttt{quydinh}, \texttt{kehoach}, \texttt{stsv}. Classifier
dùng BGE-M3 embedding làm feature cho mô hình học máy trong repo.

\texttt{QueryRouter} bổ sung cơ chế two-pass: nếu câu hỏi ngắn và confidence
thấp, router có thể dùng thêm lịch sử hội thoại để phân loại; câu hỏi dài và tự
chứa không nên bị lịch sử làm lệch chủ đề.

\texttt{CollectionSelector} ánh xạ domain sang collection:

\begin{verbatim}
ctdt    -> [ctdt]
quydinh -> [quydinh, stsv]
kehoach -> [kehoach]
stsv    -> [stsv, quydinh]
\end{verbatim}

Khi confidence thấp, selector giữ collection của domain dự đoán trước rồi mở
rộng với fallback set. Với truy vấn freshness đã route về \texttt{kehoach}, hệ
thống có logic khóa route để kế hoạch/thông báo mới không bị chìm trong tài liệu
quy định hoặc CTDT.

\subsection{Query reflection và guardrail}

\texttt{QueryReflector} nhận câu hỏi hiện tại, lịch sử và
\texttt{user\_context}, sau đó trả về truy vấn đã rewrite, prompt và entities.
Các bước chính:

\begin{enumerate}
    \item strip PII/noise không cần cho retrieval;
    \item merge profile context khi câu hỏi có tín hiệu phụ thuộc hồ sơ;
    \item gọi LLM rewrite để tạo truy vấn standalone;
    \item áp guardrail chống unresolved reference và context bleed;
    \item rewrite deterministic cho follow-up so sánh ngắn;
    \item extract entities bằng regex.
\end{enumerate}

Entities quan trọng gồm \texttt{major\_code}, \texttt{major\_name},
\texttt{cohort}, \texttt{year\_of\_study}, \texttt{course\_code},
\texttt{semester} và \texttt{academic\_year}. Guardrail đảm bảo câu hỏi mới nhất
hoặc câu hỏi chung không tự động thừa kế ngành, khóa hoặc học kỳ từ lịch sử nếu
câu hỏi hiện tại không có tín hiệu rõ ràng.

\subsection{Query decomposition}

\texttt{QueryDecomposer} dùng LLM để tách câu hỏi nhiều nguồn thành tối đa vài
subquery có collection rõ ràng:

\begin{verbatim}
{
  "subqueries": [
    {"query": "...", "collection": "ctdt"},
    {"query": "...", "collection": "quydinh"}
  ]
}
\end{verbatim}

Nếu parsing lỗi, decomposer fallback về một subquery dùng câu hỏi gốc. Decomposition
giúp các câu hỏi so sánh hoặc multi-source tránh trộn toàn bộ mục tiêu vào một
truy vấn duy nhất.

\section{RetrievalService và truy hồi lai}

\subsection{RetrievalService dùng chung}

\texttt{RetrievalService.from\_settings(settings)} là entrypoint chuẩn của tầng
retrieval. Service này tạo và giữ:

\begin{itemize}
    \item \texttt{BGEm3Embedder};
    \item \texttt{E5MultilingualEmbedder};
    \item \texttt{MultiCollectionSearch};
    \item optional \texttt{BGEReranker};
    \item optional \texttt{TavilySearchTool}.
\end{itemize}

\texttt{RAGPipeline} tạo một instance \texttt{RetrievalService} duy nhất và inject
instance này vào agent tool adapters. Đây là contract quan trọng vì nó tránh load
lại các model nặng cho classic RAG và agent.

\subsection{Qdrant dual named vectors}

Mỗi collection trong Qdrant lưu hai named vectors:

\begin{itemize}
    \item \texttt{bge\_m3}: vector 1024 chiều từ BGE-M3;
    \item \texttt{e5}: vector 1024 chiều từ E5 multilingual.
\end{itemize}

Khi vector search, Qdrant chạy cả hai vector và fuse điểm theo trọng số nội bộ.
Kết quả vector sau đó được đưa vào tầng fusion với kết quả keyword từ
Elasticsearch.

\subsection{Elasticsearch keyword và metadata search}

Elasticsearch đảm nhiệm hai nhiệm vụ:

\begin{itemize}
    \item keyword search bằng \texttt{multi\_match} trên \texttt{text} và
    \texttt{title} để bắt exact term, tên riêng, mã môn, số điều;
    \item metadata search để lấy danh sách chunk id phù hợp với filter trước khi
    Qdrant vector search.
\end{itemize}

Cơ chế metadata search rất quan trọng vì vector database cần một tập id phù hợp
để áp \texttt{HasIdCondition}. Nếu metadata filter không match kết quả nào,
fallback chain sẽ nới lỏng filter thay vì trả rỗng ngay.

\subsection{MultiCollectionSearch}

\texttt{MultiCollectionSearch} là engine tìm kiếm toàn cục. Luồng chính:

\begin{enumerate}
    \item parse structured query và exclusion terms;
    \item xác định fusion weight thích ứng;
    \item tạo filter spec theo collection;
    \item resolve filter bằng Elasticsearch metadata search;
    \item chạy vector và keyword search song song trên các collection;
    \item prefix id dạng \texttt{\{collection\}/\{id\}} để tránh trùng id;
    \item tạo global vector pool và global keyword pool;
    \item hợp nhất điểm bằng linear score fusion hoặc RRF;
    \item cộng recency bonus cho \texttt{kehoach};
    \item deduplicate theo id và text;
    \item trả top candidate cho reranker.
\end{enumerate}

Ở runtime, settings mặc định dùng \texttt{vector\_weight=0.8},
\texttt{keyword\_weight=0.2}, \texttt{vector\_top\_k=50},
\texttt{keyword\_top\_k=50}, \texttt{vector\_pool\_k=40},
\texttt{keyword\_pool\_k=40} và \texttt{top\_k=5}. Với truy vấn giống câu hỏi học
phần, hệ thống tăng trọng số keyword lên ít nhất 0.6 để ưu tiên mã môn và tên
học phần.

\subsection{Metadata filtering theo collection}

\begin{table}[h]
\centering
\begin{tabular}{|l|p{9.4cm}|}
\hline
\textbf{Collection} & \textbf{Chiến lược filter} \\
\hline
\texttt{ctdt} & Ưu tiên \texttt{major\_code} exact match, fallback sang
\texttt{major\_name} fuzzy, sau đó bao gồm chunk generic hoặc bỏ filter. \\
\hline
\texttt{quydinh} & Dùng \texttt{applicable\_cohort} từ query hoặc resolved
cohort; fallback bao gồm văn bản generic không gắn khóa. \\
\hline
\texttt{kehoach} & Nếu có tháng/năm đăng bài thì wildcard trên
\texttt{date\_str}; nếu có freshness intent thì lấy chunk mới nhất theo
\texttt{date\_str}. Token học kỳ như \texttt{2025.2} không bị coi là ngày đăng. \\
\hline
\texttt{stsv} & Không có prefilter mặc định; tìm trực tiếp theo nội dung và
metadata chung. \\
\hline
\end{tabular}
\caption{Metadata filtering theo collection}
\end{table}

\subsection{Reranking và hậu xử lý}

Sau hybrid retrieval, BGE reranker đánh điểm cặp query--document và sắp xếp lại
candidate. Ngưỡng mặc định cho chunk thường là \texttt{0.0}; chunk có bảng dùng
ngưỡng riêng \texttt{reranker\_table\_score\_threshold}, hiện mặc định
\texttt{-1.0}, để giữ bảng liên quan nhưng loại bảng sai ngành quá nhiễu.

Sau rerank, pipeline áp các bước hậu xử lý:

\begin{itemize}
    \item \textbf{ValidityFilter}: loại văn bản superseded khi còn đủ nguồn;
    \item \textbf{ReferenceResolver}: bổ sung chunk cùng văn bản khi đoạn truy
    hồi có dẫn chiếu như ``Điều 5'', ``khoản 1 Điều 5'';
    \item \textbf{context formatting}: thêm metadata header, source và profile
    note khi câu hỏi thật sự cần cá nhân hóa.
\end{itemize}

\section{Agentic RAG}

\subsection{Vai trò của agent}

Classic RAG tuyến tính phù hợp với câu hỏi một mục tiêu. Tuy nhiên, các câu hỏi
như ``so sánh chương trình IT-E6 và IT-E7'', ``điều kiện tốt nghiệp của khóa K70
khác K65 thế nào'' hoặc ``vừa hỏi quy định vừa hỏi lịch đăng ký'' cần nhiều bước
truy hồi và tổng hợp. Vì vậy hệ thống bổ sung agent trong \texttt{agent/}.

\subsection{Hai đường thực thi}

\texttt{ReActAgent} hiện có hai đường thực thi:

\begin{itemize}
    \item \textbf{Planner-executor}: dùng cho subtype \texttt{comparison} và
    \texttt{multi\_source}. Agent decompose câu hỏi, tạo retrieval plan, validate
    plan, execute song song và synthesize câu trả lời.
    \item \textbf{ReAct tool loop}: dùng cho câu hỏi phức tạp tổng quát hoặc khi
    planner fallback. Local tool-calling LLM chọn công cụ, nhận kết quả, lặp lại
    hoặc synthesize.
\end{itemize}

\begin{verbatim}
RAGPipeline.query_agent()
  -> init_agent_docs()
  -> ReActAgent.run()
     -> planner path: decompose -> planner -> validate -> executor -> synthesize
     -> react path: LLM tool call -> tools -> continue/synthesize/clarify
  -> get_agent_docs()
  -> response mapper + Mongo trace
\end{verbatim}

\subsection{Tool contract}

Các công cụ được bind trực tiếp cho ReAct LLM:

\begin{table}[h]
\centering
\begin{tabular}{|l|p{8.8cm}|}
\hline
\textbf{Tool} & \textbf{Chức năng} \\
\hline
\texttt{rag\_search} & Tìm kiếm một logical collection của RAG. \\
\hline
\texttt{web\_search} & Tìm kiếm Tavily trên domain HUST/giáo dục khi cần dữ
liệu mới hoặc thiếu nguồn nội bộ. \\
\hline
\texttt{clarify\_question} & Trả câu hỏi làm rõ và dừng lượt hiện tại. \\
\hline
\end{tabular}
\caption{Tool visible với LangGraph agent}
\end{table}

Adapter vẫn hỗ trợ các tool cũ như \texttt{multi\_rag\_search},
\texttt{compare\_cohorts} và \texttt{compare\_programs} cho test và backward
compatibility, nhưng các tool này không được bind trực tiếp cho LLM. So sánh và
multi-source nên đi qua planner-executor.

\subsection{Agent collection alias}

Agent dùng alias thân thiện hơn với domain nội bộ:

\begin{verbatim}
chuong_trinh -> ctdt
quy_dinh     -> quydinh
ke_hoach     -> kehoach
ho_tro_sv    -> stsv
\end{verbatim}

Tool adapter map alias về collection thật, strip PII/student id, dùng query đã
chuẩn hóa cho vector và giữ query phù hợp cho Elasticsearch. Reranker trong agent
được bảo vệ bằng lock vì tokenizer/runtime path không được giả định thread-safe.

\section{Web fallback và tự đánh giá}

\subsection{Tavily fallback có kiểm soát}

Tavily là thành phần optional. \texttt{RetrievalService} chỉ tạo
\texttt{TavilySearchTool} khi API key hợp lệ. Web search được giới hạn trong các
domain HUST chính thức, domain HUST mở rộng và các nguồn giáo dục có thẩm quyền.

Trong non-streaming classic RAG, có hai giai đoạn web:

\begin{itemize}
    \item \textbf{Pre-generation enrichment}: bổ sung context web khi local source
    rỗng, truy vấn dynamic/freshness hoặc retrieval confidence thấp;
    \item \textbf{Post-generation fallback}: regenerate khi câu trả lời báo
    không có thông tin, không có source hoặc self-eval yêu cầu web search với
    trạng thái \texttt{insufficient}/\texttt{stale\_risk}.
\end{itemize}

Cả hai giai đoạn đều phụ thuộc \texttt{tavily\_fallback\_enabled}. Streaming path
không chạy post-generation Tavily để không phá vỡ trải nghiệm stream.

\subsection{SelfEvaluator}

Self-evaluation được dùng như quality gate tùy chọn. Nếu direct self-eval tắt
nhưng Tavily fallback bật, pipeline vẫn có thể khởi tạo evaluator vì self-eval
là tín hiệu quyết định fallback web. Ngưỡng
\texttt{self\_eval\_min\_top\_score} mặc định rất cao vì BGE reranker trả raw
logit, không phải xác suất.

\section{API, xác thực và contract response}

\subsection{Chat API}

Các endpoint chat chính:

\begin{itemize}
    \item \texttt{POST /chat}: gọi \texttt{query\_v3()} và map về
    \texttt{ChatResponse};
    \item \texttt{POST /chat/v3}: response giàu trace/debug cho web/mobile;
    \item \texttt{POST /api/chat/v3}: alias của \texttt{/chat/v3};
    \item \texttt{POST /chat/stream}: SSE streaming;
    \item \texttt{GET /chat/suggest}: câu hỏi gợi ý cho mobile/profile context.
\end{itemize}

Nếu có Bearer JWT hợp lệ, route lấy \texttt{user\_id} và
\texttt{user\_context} từ DB user, không tin identity spoofable trong request
body. Request body vẫn giữ các field legacy để hỗ trợ client cũ hoặc anonymous
mode.

\subsection{Response metadata}

\texttt{api/response\_mapper.py} chuẩn hóa output không đồng nhất của classic
RAG và agent thành schema cho client. Metadata có thể gồm:

\begin{itemize}
    \item answer và retrieved documents;
    \item target collections, collection scores, collection results;
    \item reflected question, applied filters, routing probabilities;
    \item mode, route, tools, tool calls, iterations;
    \item agent trace, agent error;
    \item timings, session id, turn id;
    \item prompt/debug fields khi có.
\end{itemize}

Nhờ response mapper, backend có thể thay đổi chi tiết nội bộ của pipeline nhưng
vẫn giữ contract tương đối ổn định cho web/mobile.

\subsection{Auth và RBAC}

Auth route nằm ở \texttt{routers/auth.py} với prefix \texttt{/auth}. JWT chứa
\texttt{sub}, \texttt{email}, \texttt{role}, \texttt{iat}, \texttt{exp}. Các
admin endpoint dùng \texttt{require\_admin}; superadmin là overlay cấu hình qua
\texttt{SUPERADMIN\_USER\_IDS}, không phải role DB riêng.

Microsoft OAuth được hỗ trợ cho tài khoản HUST/SIS. User document lưu email,
username, role, student id, full name, cohort, major và \texttt{major\_code} để
phục vụ context hội thoại.

\subsection{API quản trị, lookup và mobile features}

Ngoài chat, backend cung cấp:

\begin{itemize}
    \item \texttt{/admin/documents*}: upload, convert, clean, chunk, approve,
    index, rollback tài liệu;
    \item \texttt{/retrieval/search}: endpoint diagnostic cho raw retrieval;
    \item \texttt{/metrics/usage}, \texttt{/metrics/eval}: metrics và eval
    dashboard;
    \item \texttt{/session}, \texttt{/sessions}, \texttt{/sessions/me}: quản lý
    phiên hội thoại;
    \item \texttt{/lookup/ctdt/\{major\_code\}},
    \texttt{/lookup/regulations}, \texttt{/lookup/calendar},
    \texttt{/lookup/compare}: tra cứu nhanh cho mobile;
    \item \texttt{/bookmarks}, \texttt{/bookmark-folders}, \texttt{/feedback},
    \texttt{/notifications}: các tính năng mobile/web phụ trợ.
\end{itemize}

\section{Persistence, cache và quan sát}

\subsection{MongoDB}

MongoDB là tầng lưu trữ bền vững. Các collection quan trọng:

\begin{table}[h]
\centering
\begin{tabular}{|l|p{9.0cm}|}
\hline
\textbf{Collection} & \textbf{Vai trò} \\
\hline
\texttt{users} & Tài khoản, role, profile và context sinh viên. \\
\hline
\texttt{sessions} & Metadata phiên chat. \\
\hline
\texttt{turns} & Lượt hỏi đáp, answer, sources và metadata. \\
\hline
\texttt{query\_logs} & Analytics log phẳng theo turn. \\
\hline
\texttt{agent\_traces} & Trace LangGraph/tool calls của agent. \\
\hline
\texttt{documents}, \texttt{document\_chunks} & Tài liệu admin upload và chunks
review/indexing. \\
\hline
\texttt{system\_config} & Override LLM/admin config, đặc biệt document
\texttt{llm\_config}. \\
\hline
\end{tabular}
\caption{Các MongoDB collection chính}
\end{table}

\subsection{Redis}

Redis là thành phần fail-soft, điều khiển bởi \texttt{redis\_enabled} và các flag
con. Khi Redis không khả dụng, backend tiếp tục qua MongoDB hoặc bỏ qua cache.

Các nhóm key chính:

\begin{itemize}
    \item \texttt{session:\{sid\}} và \texttt{user\_sessions:\{uid\}} cho session
    metadata/list;
    \item \texttt{history:\{sid\}} cho context hội thoại gần đây;
    \item \texttt{llm\_cache:\{sha\}} và \texttt{llm\_cache:q:\{sha\}} cho cache
    post-retrieval và query-only;
    \item \texttt{rate:min:\{id\}}, \texttt{rate:day:\{id\}} cho rate limiting.
\end{itemize}

Rate limit mặc định là \texttt{20} request/phút và \texttt{200} request/ngày.
Nếu Redis lỗi, middleware cho request đi qua và log warning thay vì làm hệ thống
fail-closed.

\subsection{Trace và observability}

Trace metadata được ghi ở nhiều tầng:

\begin{itemize}
    \item route/domain confidence và routing probabilities;
    \item reflected question và extracted entities;
    \item filter đã áp dụng, số id match và collection counts;
    \item fusion weights, structured query, excluded counts;
    \item rerank scores, validity/reference resolver behavior;
    \item timings từng bước;
    \item agent tool calls, iterations, route và errors.
\end{itemize}

Thông tin này phục vụ debug trong frontend trace page, phân tích lỗi retrieval và
đánh giá chất lượng sau khi cập nhật dữ liệu.

\section{Giao diện web, mobile và shared package}

\subsection{Web frontend}

Web app nằm ở \texttt{frontend/chat-companion}. Các route chính gồm
\texttt{/chat}, \texttt{/chat/:sessionId}, \texttt{/login}, \texttt{/register},
\texttt{/complete-profile}, \texttt{/trace}, \texttt{/retrieval},
\texttt{/eval}, \texttt{/admin}, \texttt{/admin/documents/:id},
\texttt{/bookmarks} và \texttt{/notifications}.

\texttt{ChatContainer} điều phối gửi tin nhắn, streaming token, session history,
metadata/debug panel và invalidation. Admin UI dùng \texttt{services/adminApi.ts}
để upload, review, chunk và index tài liệu. Trace UI đọc metadata từ
\texttt{/chat/v3} hoặc stream metadata để hiển thị pipeline.

\subsection{Mobile app}

Mobile app nằm ở \texttt{mobile/} và dùng Expo/React Native. Main tabs gồm Chat,
Lookup, Bookmarks, Notifications và Profile. Token được lưu trong SecureStore;
cache offline không nhạy cảm dùng MMKV khi native runtime hỗ trợ, còn Expo Go có
fallback in-memory.

Streaming mobile dùng \texttt{/chat/stream} qua \texttt{react-native-sse}; nếu
streaming lỗi trước token đầu tiên, client có thể fallback sang \texttt{/chat/v3}.

\subsection{Shared TypeScript package}

\texttt{packages/shared} cung cấp API clients, types, stores và utilities dùng
chung. Các path như \texttt{/chat/v3}, \texttt{/chat/stream},
\texttt{/auth/login}, \texttt{/sessions/me}, \texttt{/lookup/calendar},
\texttt{/notifications} được khai báo tập trung để web/mobile giảm lệch contract.

\section{Đánh giá và hồi quy}

\subsection{Hai lớp đánh giá}

Module \texttt{evaluation/} tách đánh giá thành hai lớp:

\begin{itemize}
    \item \textbf{Current policy eval}: đo retrieval theo dữ liệu và chính sách
    hiện hành, gồm collection accuracy, recall, NDCG, MRR, context
    precision/recall, citation accuracy và freshness;
    \item \textbf{Historical conversation eval}: đo hiểu hội thoại, follow-up,
    personalization, clarification và advisory behavior trên dữ liệu lịch sử.
\end{itemize}

Sự tách biệt này cần thiết vì câu trả lời lịch sử có thể khác chính sách hiện
hành. Mismatch với email cũ không nên tự động coi là lỗi factual nếu tài liệu
hiện tại đã thay đổi.

\subsection{Các check phù hợp sau thay đổi}

Theo cấu trúc module hiện tại, các thay đổi cần chọn test tương ứng:

\begin{itemize}
    \item thay đổi query/routing/reflection: chạy test query, phase7/phase8 và
    regression conversation;
    \item thay đổi retrieval/fusion/filter: chạy \texttt{tests/retrieval},
    \texttt{test\_multi\_collection\_fusion.py} và current policy eval;
    \item thay đổi document pipeline: chạy \texttt{tests/test\_document\_pipeline.py}
    và \texttt{tests/test\_upload\_api.py};
    \item thay đổi response metadata: cập nhật mapper, web/mobile normalizer và
    trace components;
    \item sau indexing/crawling lớn: chạy post-index hoặc current policy eval.
\end{itemize}

\section{Tổng hợp điểm mới của phương pháp đề xuất}

Phương pháp đề xuất trong hệ thống \texttt{RAG\_v2} có các điểm nổi bật:

\begin{itemize}
    \item \textbf{Smart entrypoint}: \texttt{query\_v3()} chọn chitchat, clarify,
    classic RAG, decomposed RAG hoặc agent theo độ phức tạp câu hỏi.
    \item \textbf{RetrievalService dùng chung}: classic RAG và agent dùng cùng
    embedders, searcher, reranker và Tavily tool để tránh load model trùng.
    \item \textbf{Hybrid retrieval đa collection}: Qdrant dual vectors kết hợp
    Elasticsearch BM25, metadata prefilter, global fusion, adaptive weights và
    recency bonus.
    \item \textbf{Guardrail hội thoại}: reflection chỉ dùng profile/history khi
    câu hỏi có tín hiệu cần cá nhân hóa, giảm context bleed.
    \item \textbf{Validity và reference resolution}: xử lý văn bản superseded và
    dẫn chiếu điều/khoản trong cùng tài liệu.
    \item \textbf{Agentic path}: planner-executor cho so sánh/nhiều nguồn và
    ReAct tool loop cho câu hỏi phức tạp tổng quát.
    \item \textbf{Admin document lifecycle}: tài liệu mới có thể được upload,
    review, chunk, index và rollback qua API/UI quản trị.
    \item \textbf{Observability}: route, filter, scores, timings, agent traces và
    eval artifacts giúp phân tích lỗi và cải tiến hệ thống.
\end{itemize}

\section{Kết luận chương}

Chương này đã mô tả đầy đủ phương pháp đề xuất của hệ thống \texttt{RAG\_v2}
dựa trên source code hiện tại. Kiến trúc được tổ chức quanh \texttt{RAGPipeline}
và \texttt{RetrievalService}, kết hợp query understanding, hybrid retrieval,
reranking, validity filtering, Agentic RAG, web fallback, persistence, cache,
admin document pipeline và client web/mobile.

Thiết kế này cho phép hệ thống trả lời câu hỏi học vụ theo hướng có căn cứ,
truy vết được nguồn, mở rộng được dữ liệu và xử lý được cả truy vấn đơn giản lẫn
truy vấn nhiều bước. Các chương tiếp theo có thể đi sâu vào triển khai từng
pipeline, thuật toán cụ thể và kết quả thực nghiệm.

\end{document}
